% Options for packages loaded elsewhere
\PassOptionsToPackage{unicode}{hyperref}
\PassOptionsToPackage{hyphens}{url}
%
\documentclass[
  brazilian,
]{article}
\usepackage{lmodern}
\usepackage{amssymb,amsmath}
\usepackage{ifxetex,ifluatex}
\ifnum 0\ifxetex 1\fi\ifluatex 1\fi=0 % if pdftex
  \usepackage[T1]{fontenc}
  \usepackage[utf8]{inputenc}
  \usepackage{textcomp} % provide euro and other symbols
\else % if luatex or xetex
  \usepackage{unicode-math}
  \defaultfontfeatures{Scale=MatchLowercase}
  \defaultfontfeatures[\rmfamily]{Ligatures=TeX,Scale=1}
\fi
% Use upquote if available, for straight quotes in verbatim environments
\IfFileExists{upquote.sty}{\usepackage{upquote}}{}
\IfFileExists{microtype.sty}{% use microtype if available
  \usepackage[]{microtype}
  \UseMicrotypeSet[protrusion]{basicmath} % disable protrusion for tt fonts
}{}
\makeatletter
\@ifundefined{KOMAClassName}{% if non-KOMA class
  \IfFileExists{parskip.sty}{%
    \usepackage{parskip}
  }{% else
    \setlength{\parindent}{0pt}
    \setlength{\parskip}{6pt plus 2pt minus 1pt}}
}{% if KOMA class
  \KOMAoptions{parskip=half}}
\makeatother
\usepackage{xcolor}
\IfFileExists{xurl.sty}{\usepackage{xurl}}{} % add URL line breaks if available
\IfFileExists{bookmark.sty}{\usepackage{bookmark}}{\usepackage{hyperref}}
\hypersetup{
  pdftitle={Parte III: Processamento de Linguagem Natural para Linguistas},
  pdfauthor={CiberExt 26-29 · FEELT38103 · Universidade Federal de Uberlândia},
  pdflang={pt-BR},
  hidelinks,
  pdfcreator={LaTeX via pandoc}}
\urlstyle{same} % disable monospaced font for URLs
\setlength{\emergencystretch}{3em} % prevent overfull lines
\providecommand{\tightlist}{%
  \setlength{\itemsep}{0pt}\setlength{\parskip}{0pt}}
\setcounter{secnumdepth}{-\maxdimen} % remove section numbering
\ifxetex
  % Load polyglossia as late as possible: uses bidi with RTL langages (e.g. Hebrew, Arabic)
  \usepackage{polyglossia}
  \setmainlanguage[]{brazilian}
\else
  \usepackage[shorthands=off,main=brazilian]{babel}
\fi

\title{Parte III: Processamento de Linguagem Natural para Linguistas}
\usepackage{etoolbox}
\makeatletter
\providecommand{\subtitle}[1]{% add subtitle to \maketitle
  \apptocmd{\@title}{\par {\large #1 \par}}{}{}
}
\makeatother
\subtitle{Unidade 2 - Exercícios}
\author{CiberExt 26-29 · FEELT38103 · Universidade Federal de
Uberlândia}
\date{Agosto de 2026}

\begin{document}
\maketitle

\hypertarget{exercuxedcios-pruxe1ticos-processamento-de-linguagem-natural-para-linguistas}{%
\section{Exercícios Práticos: Processamento de Linguagem Natural para
Linguistas}\label{exercuxedcios-pruxe1ticos-processamento-de-linguagem-natural-para-linguistas}}

Este documento contém exercícios práticos baseados na unidade
``Dependências Universais''.

\hypertarget{muxf3dulo-2-dependuxeancias-universais}{%
\subsection{Módulo 2: Dependências
Universais}\label{muxf3dulo-2-dependuxeancias-universais}}

\textbf{Exercício 1:} O projeto \emph{Universal Dependencies} (UD) é
descrito como um esforço colaborativo que possui dois objetivos
principais e sobrepostos. Quais são esses dois objetivos?

\textbf{Exercício 2:} A descrição de estruturas linguísticas dentro do
framework UD baseia-se em três tipos principais de unidades frasais.
Liste quais são essas três unidades e descreva brevemente o propósito ou
o núcleo (a que classe de palavras estão associadas) de cada uma.

\textbf{Exercício 3:} Na árvore de análise sintática (parse tree) gerada
pelo framework UD para uma oração completa, o que significa a etiqueta
de dependência \texttt{root} e como as relações de dependência sintática
fluem a partir dela em relação aos nominais da frase?

\textbf{Exercício 4:} Ao explorar as dependências sintáticas utilizando
os objetos \emph{Token} da biblioteca spaCy, você pode navegar na árvore
de análise de várias formas. Qual é a diferença exata entre os
resultados retornados pelos atributos \texttt{token.subtree} e
\texttt{token.children}?

\textbf{Exercício 5:} Por que as métricas tradicionais de avaliação,
como UAS (\emph{Unlabeled Attachment Score}) e LAS (\emph{Labeled
Attachment Score}), tornam-se problemáticas e desiguais ao comparar o
desempenho de modelos analíticos entre diferentes idiomas,
especificamente comparando o inglês com um idioma morfologicamente rico
como o finlandês? Cite ao menos uma métrica alternativa sugerida para
resolver esse problema.

\begin{center}\rule{0.5\linewidth}{0.5pt}\end{center}

\end{document}
